\section{The proximal stage of Section~\ref{sec:diagcert}}\label{app:proximal}


The following generator finds a point near $S$ when a guessed ratio
$K\ge1$ satisfies $K\ge L/g$. Its work is bounded for every guess, but its
conclusion is not checkable, so it is used only to propose candidates.

Given a rational $K\ge1$, let $\mu\ge2$ be the least integer with
$4^{-\mu}K\le1/8$, and put
\[
\theta=2^{-\mu},\qquad\eta=\frac{L\theta^2}4,\qquad
\varrho=2\bigl\lceil\sqrt{Kn}\,\bigr\rceil,\qquad
K_\theta=100\,\theta^{-1}\lceil\log_2(n+2)\rceil .
\]
Then $\theta\le1/4$, $\theta^2K\le1/8$ and $\theta^{-1}\le6\sqrt K$.
A \emph{proximal stage} with center $c\in X$ and mesh $h>0$ does the
following. Form the fresh box $B=X\cap\prod_i[c_i-\varrho h,c_i+\varrho h]$,
rounding integer endpoints inward. Build in each $B_i$ the outward geometric
grid $G_i$ from $c_i$ of Section~\ref{sec:growth}, with step $h+\theta t$
(continuous) or $\max\{1,\lfloor h+\theta t\rfloor\}$ (integer) at distance
$t$, clipped at the endpoints of $B_i$. With the corrections $d_i$ and $D$ of
\eqref{eq:correction} on these grids, minimize
\[
P(y)=F(y)-D(y)+\eta\norm{y-c}^2\qquad\Bigl(y\in\prod_iG_i\Bigr)
\]
by the junction-tree recursion \eqref{eq:messages}; correction and proximal
terms are unary, so the decomposition is unchanged. The output is a
minimizer $y^+$. The \emph{proximal iteration with guess $K$} starts at the
lower endpoint vector $c^0=l$, uses $h_j=s2^{-j}$, outputs $y^j$ at stage $j$,
and sets $c^{j+1}=y^j$. Every center is feasible, and every box is built from
$X$ rather than from the previous box.

\begin{lemma}[Proximal stage]\label{lem:proximal}
\begin{enumerate}
\item[(i)] For every guess, every grid has at most $K_\theta$ nodes.
\item[(ii)] Let $g$ satisfy \eqref{eq:setgrowth} with $L/g\le K$. If
$\dist(c,S)^2\le4Knh^2$, then
\[
F(y^+)-\OPT\le\tfrac{99}{256}Lnh^2,\qquad
\dist(y^+,S)^2\le\tfrac{99}{256}Knh^2 .
\]
\item[(iii)] Under the hypothesis of (ii) on $g$, the proximal iteration
satisfies $\dist(y^j,S)^2\le\frac{99}{256}Knh_j^2$ for every $j\ge0$.
\item[(iv)] Stage $j$ costs $O\bigl(p(N+|\mathcal A|)K_\theta^p\bigr)$ table
operations on rationals of bit length $\poly(I+j+\mu K_\theta)$, for every
guess.
\end{enumerate}
\end{lemma}

\begin{proof}
(i) Each endpoint of $B_i$ is within $\varrho h$ of $c_i$. As in the proof of
Lemma~\ref{lem:states}, one side of a continuous grid has at most
$1+\lceil\ln(1+\theta\varrho)/\ln(1+\theta)\rceil$ nodes, and one side of an
integer grid at most $1+\lceil\ln(1+\theta\varrho h/H)/\ln(1+\theta/3)\rceil$
with $H=\max\{h,1\}\ge h$. Now
$\theta\varrho\le2\theta\sqrt{Kn}+2\theta\le\sqrt{n/2}+\frac12$, and
$\frac32+\sqrt{n/2}\le n+2$. Using $\ln(1+\theta)\ge\ln(1+\theta/3)\ge\theta/4$
and $\ln(n+2)\le\log_2(n+2)$, each side has at most
$2+4\theta^{-1}\log_2(n+2)$ nodes. With the center, $|G_i|\le5+8\theta^{-1}
\lceil\log_2(n+2)\rceil\le10\,\theta^{-1}\lceil\log_2(n+2)\rceil\le K_\theta$,
since $\theta^{-1}\ge4$.

(ii) Let $s^*\in S$ be nearest to $c$ (it exists since $S$ is compact) and
$E=\norm{s^*-c}^2\le4Knh^2$. Each $|s^*_i-c_i|\le2\sqrt{Kn}\,h\le\varrho h$,
and integer coordinates of $s^*$ are integers, so $s^*\in B$. Let $Y$ be the
independent mean-preserving rounding of $s^*$ to its enclosing grid
intervals. Proposition~\ref{prop:cellwise} on the box $B$ gives
$\E[F(Y)-D(Y)]\le F(s^*)=\OPT$. Independence gives
$\E\norm{Y-c}^2=E+\sum_i\operatorname{Var}Y_i$. If $s^*_i$ is not a node, its
enclosing interval was generated from its inner endpoint, at distance at most
$|s^*_i-c_i|$ from $c_i$, so its length is at most $h+\theta|s^*_i-c_i|$;
an integer coordinate inside a unit interval is impossible, and clipping only
shortens. Hence
$\operatorname{Var}Y_i\le\frac14(h+\theta|s^*_i-c_i|)^2\le\frac12(h^2+\theta^2|s^*_i-c_i|^2)$
and $\sum_i\operatorname{Var}Y_i\le\frac12(nh^2+\theta^2E)$. Since a minimum
is at most an expectation,
\[
P(y^+)\le\OPT+\eta\Bigl[\bigl(1+\tfrac{\theta^2}2\bigr)E+\tfrac{nh^2}2\Bigr].
\]
By Lemma~\ref{lem:graded}(a), $w_i(v)\le h+\theta|v-c_i|$ for every node, so
$d_i(v)\le\frac L4h^2+\eta|v-c_i|^2$ and
$D(y^+)\le\frac L4nh^2+\eta\norm{y^+-c}^2$. Therefore
\begin{align*}
F(y^+)-\OPT&=P(y^+)-\OPT+D(y^+)-\eta\norm{y^+-c}^2\\
&\le\frac{Lnh^2}4\Bigl[1+4K\theta^2\bigl(1+\tfrac{\theta^2}2\bigr)+\tfrac{\theta^2}2\Bigr]
\le\frac{Lnh^2}4\cdot\frac{99}{64},
\end{align*}
using $E\le4Knh^2$, $4K\theta^2\le\frac12$ and $\theta^2\le\frac1{16}$.
Growth gives $\dist(y^+,S)^2\le(F(y^+)-\OPT)/g\le\frac{99}{256}(L/g)nh^2\le
\frac{99}{256}Knh^2$.

(iii) Initially $\dist(l,S)^2\le ns^2\le4Knh_0^2$. If
$\dist(y^j,S)^2\le\frac{99}{256}Knh_j^2$, then
$\dist(c^{j+1},S)^2\le Knh_j^2/2=2Knh_{j+1}^2\le4Knh_{j+1}^2$. Apply (ii).

(iv) Let $\omega_0$ be a common denominator of $H,b,c,L$ and the endpoints.
By induction, every stage-$j$ grid coordinate has a denominator dividing
$\omega_02^{j+\mu K_\theta}$: the center comes from stage $j-1$; the new
endpoints $c_i\pm\varrho h_j$ have denominators dividing $\omega_02^j$; an
unclipped continuous offset is
$h_j\bigl[(2^\mu+1)^k-2^{\mu k}\bigr]/2^{\mu(k-1)}$ with $k\le K_\theta$;
integer steps are integral; clipping selects an endpoint of $B_i$. Values of
$F$, corrections and the proximal term then have a common denominator
dividing $8\omega_0^22^{2\mu+2}\bigl(\omega_02^{j+\mu K_\theta}\bigr)^2$, and
bounded magnitudes because all points stay in $\bar X$. Messages are sums of
such values. The table count is that of Lemma~\ref{lem:dp} with $K_\theta$ labels.
\end{proof}

\begin{remark}[A false guess]\label{rem:falseguess}
Without $L/g\le K$ the iteration may converge to a nonglobal point. For
$F=x^2+y^2-3xy+\frac{63}{128}(x+y)$ on $[0,1]^2$, put $t=x+y$; then
$F\ge t^2-\frac54t^2+\frac{63}{128}t$, a concave function of $t\in[0,2]$, so
$(1,1)$ is the unique optimizer with $\OPT=-\frac1{64}$. The origin is a box
KKT point with $F=0$. Exact computation shows that the iteration with
$K\in\{1,2,4\}$ returns the origin at every stage $j\le33$, which covers the
budgets $J_K$ of Theorem~\ref{thm:diagdiscovery}. Lemma~\ref{lem:diagcert}
rejects it: $H+\Lambda=\bigl(\begin{smallmatrix}191/64&-3\\-3&191/64\end{smallmatrix}\bigr)$
has eigenvalue $-\frac1{64}$. At $(1,1)$ the matrix has diagonal $\frac{193}{64}$
and is positive definite.
\end{remark}

